%%%%%%%%%%%%%%%%%%%%%%%%%%%%%%%%%%%%%%%%%%%%%%%%%%%%%%%%%%%%%%%%%%%%%%%%%%%%
% CAPÍTULO 6: ALCANCE Y LIMITACIONES
%%%%%%%%%%%%%%%%%%%%%%%%%%%%%%%%%%%%%%%%%%%%%%%%%%%%%%%%%%%%%%%%%%%%%%%%%%%%

\chapter{Alcance y limitaciones}
\label{ch:alcance}

\section{Alcance del proyecto}

\subsection{Dentro del alcance}

Los elementos que se incluyen dentro del alcance del proyecto se detallan en la Tabla~\ref{tab:dentro-alcance}.

\begin{table}[H]
  \centering
  \begin{tabular}{@{}p{\textwidth}@{}}
    \hline
    \textbf{Elemento} \\
    \hline
    Registro con verificación del correo, inicio de sesión único y recuperación de la contraseña \\
    Perfiles de conductor, pasajero y administrador (selector de modo) \\
    Selección de direcciones y de la ubicación del dispositivo, y los parkings del campus como destino \\
    Publicación y solicitud de viajes \\
    \textit{Matching} conductor-pasajero \\
    Gestión de plazas y estado del viaje \\
    Representación mediante Gemelo Digital \\
    Simulación de escenarios \\
    Métricas y evaluación experimental \\
    CI/CD y documentación técnica \\
    \hline
  \end{tabular}
  \caption{Elementos dentro del alcance.}
  \label{tab:dentro-alcance}
\end{table}

El ámbito geográfico se limita a una \textbf{zona de servicio}: un círculo de 30~km de radio alrededor del campus de La Cañada, cuyo centro es la media de la posición de sus parkings. Abarca Almería y los municipios cercanos desde los que se viaja a diario al campus (como Huércal de Almería, Vícar o Roquetas de Mar), y deja fuera los más alejados de la provincia. Una dirección fuera de la zona no puede usarse para buscar ni publicar viajes (RF-27): la aplicación avisa de que aún no hay cobertura allí y registra, de forma anónima, el municipio y una zona aproximada de 1~km, para conocer la demanda no cubierta sin guardar datos personales (RNF-04). La zona se define en la base de datos, de modo que se puede ampliar o ajustar sin cambiar el código.

\subsection{Fuera del alcance inicial}

Los elementos que quedan fuera del alcance inicial y sus motivos se muestran en la Tabla~\ref{tab:fuera-alcance}.

\begin{table}[H]
  \centering
  \begin{tabularx}{\textwidth}{Y{0.7692} Y{1.2308}}
    \hline
    \textbf{Elemento} & \textbf{Motivo} \\
    \hline
    Integración obligatoria con sistemas institucionales de identidad de la UAL & Dependencia externa no controlada \\
    Despliegue productivo para toda la universidad & Escala no requerida para demostrar la contribución \\
    Pagos reales entre usuarios & Fuera del foco funcional del TFM \\
    Gestión financiera o facturación & No aporta a la pregunta de investigación \\
    Vehículos autónomos & Modelo operativo distinto (B5, B6) \\
    Control físico de vehículos & Requiere hardware no disponible \\
    Predicción avanzada obligatoria con modelos de IA complejos & No necesaria para el MVP \\
    Optimización multiagente de tráfico urbano a escala de ciudad & Excede el alcance del prototipo \\
    Infraestructura 6G/\textit{edge} especializada & No aporta al problema estudiado (E11) \\
    Cálculo de métricas en tiempo real & Se calculan al finalizar cada simulación \\
    \hline
  \end{tabularx}
  \caption{Elementos fuera del alcance inicial.}
  \label{tab:fuera-alcance}
\end{table}

\subsection{Funcionalidades recortables}
\label{sec:moscow-alcance}

Estas funcionalidades no forman parte del MVP, pero se implementarán si el avance del proyecto lo permite. Mientras no se implementen, se usa la alternativa indicada: avisos dentro de la aplicación (RF-19) en lugar de notificaciones \textit{push} o por \textit{email}, vista de resultados con exportación CSV (RF-20) en lugar del panel con gráficos, configuraciones predefinidas (RF-16) en lugar de la parametrización libre, reglas de aviso fijas (RF-19) en lugar de su configuración por perfil, y avisos dentro de la aplicación (RF-19) para coordinar a conductor y pasajeros en lugar del chat. En caso contrario, se pospondrán a trabajo futuro sin afectar a la validación del TFM (véase la Tabla~\ref{tab:moscow}).

\begin{table}[H]
  \centering
  \begin{tabularx}{\textwidth}{@{}Y{1.0268} Y{0.4027} Y{1.5705}@{}}
    \hline
    \textbf{Elemento} & \textbf{Prioridad} & \textbf{Condición} \\
    \hline
    Notificaciones \textit{push} (FCM) & Debería & Se implementa si F2--F5 se cierran en plazo \\
    Notificaciones por \textit{email} (SMTP) & Debería & Se implementa si F2--F5 se cierran en plazo \\
    Panel de administrador con gráficos interactivos & Podría & Se implementa si sobra tiempo tras el MVP \\
    Parametrización libre de escenarios & Podría & Se implementa si sobra tiempo tras el MVP \\
    Configuración de avisos por perfil (cada modo configura los suyos) & Podría & Se implementa si sobra tiempo tras el MVP \\
    Chat entre el conductor y los pasajeros aceptados de un viaje & Podría & Se implementa si sobra tiempo tras el MVP \\
    Límite de intentos de inicio de sesión por cuenta e IP & Podría & Se implementa si sobra tiempo tras el MVP \\
    Segundo factor (TOTP) para los administradores & Podría & Se implementa si sobra tiempo tras el MVP \\
    Verificación del teléfono por SMS & Podría & Requiere un proveedor de pago; hoy solo se valida el formato \\
    Aplicación instalable (PWA) & Podría & Se implementa si sobra tiempo tras el MVP \\
    Idiomas español e inglés & Podría & Pensado para estudiantes Erasmus y profesorado visitante \\
    \hline
  \end{tabularx}
  \caption{Funcionalidades recortables (MoSCoW).}
  \label{tab:moscow}
\end{table}

\subsection{Criterio de alcance}

El Gemelo Digital debe aportar valor al problema del coche compartido. No se pretende construir una réplica digital de toda la ciudad ni resolver el transporte urbano completo, sino modelar el sistema específico de viajes compartidos que gestiona la plataforma. Este criterio actúa como filtro para las decisiones de diseño: si un componente no contribuye a representar, gestionar o evaluar el coche compartido universitario, queda fuera del alcance.

\section{Limitaciones previstas}

Las limitaciones se separan de las exclusiones de alcance porque no son decisiones de diseño, sino restricciones reales del trabajo. Se detallan en la Tabla~\ref{tab:limitaciones}.

\begin{table}[H]
  \centering
  \begin{tabularx}{\textwidth}{Y{0.6923} Y{1.8462} Y{0.4615}}
    \hline
    \textbf{Limitación} & \textbf{Descripción} & \textbf{Impacto} \\
    \hline
    Datos reales & No se dispone inicialmente de datos reales de movilidad de la UAL; se usarán datos sintéticos y escenarios reproducibles. & Medio \\
    Fidelidad del Gemelo Digital & El DT es de baja fidelidad: representa estado y eventos, no dinámica física ni sensórica. & Medio \\
    Escala & El prototipo se diseña para un entorno controlado, no para despliegue masivo. & Bajo \\
    \textit{Matching} & La primera versión usa \textit{matching} heurístico interpretable, no optimización avanzada. & Bajo \\
    Validación & Los resultados de simulación representan el comportamiento del modelo bajo condiciones definidas, no predicciones del mundo real. & Medio \\
    Acceso bibliográfico & Algunos trabajos relevantes quedaron pendientes de evaluación por falta de acceso al texto completo. & Bajo \\
    Trazabilidad de la SLR & No se conservaron individualmente todas las decisiones de descarte en las primeras fases de cribado. & Bajo \\
    Tiempo & El calendario del TFM limita la profundidad de algunas fases (por ejemplo, pruebas con usuarios reales). & Medio \\
    \hline
  \end{tabularx}
  \caption{Limitaciones previstas.}
  \label{tab:limitaciones}
\end{table}

\section{Relación con el MVP}

El MVP definido en el plan de trabajo constituye el núcleo mínimo que permite demostrar la contribución:
\begin{itemize}
    \item Usuarios con perfiles de conductor, de pasajero y de administrador.
    \item Publicación y búsqueda de viajes.
    \item \textit{Matching} básico.
    \item Gestión de solicitudes y estados.
    \item Modelo de Gemelo Digital de usuarios, vehículos, viajes, solicitudes y rutas.
    \item Al menos tres escenarios de simulación.
    \item Métricas comparables entre escenario base y escenarios de carpooling.
    \item Pruebas automatizadas básicas y \textit{pipeline} CI/CD.
\end{itemize}

Las funcionalidades de segunda prioridad (navegación avanzada, geolocalización real, IA predictiva, optimización multiobjetivo, integración institucional, despliegue a escala) quedan explícitamente fuera del MVP y se recogen como trabajo futuro.

\section{Amenazas a la validez}

Las amenazas a la validez de los resultados (interna, externa y de constructo) se analizan junto al plan de evaluación, en la Sección~\ref{sec:amenazas-validez}.
